Приведем первый алгоритм, который будет строить нам миностов -- алгоритм Крускала.

Будем строить миностов поэтапно. Отсортируем рёбра по весу по возрастанию, а далее будем пытаться жадно строить из них дерево.
Осталось только понять как строить дерево.

Изначально примем, что наш остов состоит только из вершин графа $G$ без каких-либо рёбер.
То есть получается каждая вершина на первом этапе представляет собой одну компоненту связности.
По сути в начале у нас есть $|V|$ \textit{непересекающихся множеств} $\{\{v_1\}, \{v_2\}, \dots, \{v_{|V|}\} \}$. Будем представлять их с помощью СНМ.
Далее будем итерироваться по рёбрам, пусть мы рассматриваем ребро $(v_i, v_j)$. 
Если $v_i$ и $v_j$ находятся в разных компонентах связности, то мы можем добавить данное ребро в дерево и объединить соответствующие компоненты связности в СНМ.
После итерации по всем рёбрам мы получим дерево.

\Exe Докажите, что на выходе мы действительно получим дерево.

Теперь представим код, выполняющий данный алгоритм. $\texttt{SORT\_BY\_W(}E, w\texttt{)}$ -- сортирует рёбра по весу.

\begin{algorithm}[H]
\caption{Алгоритм Крускала}
\begin{algorithmic}[1]

\Function{KRUSKAL}{$G\langle V, E \rangle, w$}
    \State $T \gets \{\}$
    \For{$v \in V$}
        \State $\texttt{MAKE\_SET(}v\texttt{)}$
    \EndFor

    \State $\texttt{SORT\_BY\_W(}E, w\texttt{)}$
    
    \For{$(u, v) \in E$}
        \If{$\texttt{!SAME(}u, v\texttt{)}$}
            \State $T \gets T \cup \{(u, v)\}$
            \State $\texttt{UNITE(}u, v\texttt{)}$
        \EndIf
    \EndFor

    \State \Return $T$
\EndFunction
\end{algorithmic}
\end{algorithm}

\St \texttt{KRUSKAL} работает за $\O(|E| \log |V|)$, если сортировка работает за $\O(|E| \log |E|)$, а также СНМ используется с весовой эвристикой и эвристикой сжатия путей.
\begin{proof}
    
    Инициализация $T$ работает за $\O(1)$, 3-4 строчки работают $\O(|V|)$. 
    Далее сортировка работает за $\O(|E| \log |E|)$. А вот цикл на 6-ой строчке работает за $\O(|E| \alpha(|V|))$.

    Значит итоговое время работы есть:
    \begin{multline*}
        \O(|E| \log |E| + |E| \alpha(|V|)) = \O(|E| (\log |E| + \alpha(|V|))) = \\ 
        = \O\left(|E| \left(\log |V|^2 + \alpha(|V|)\right)\right) = \O\left(|E| \left(\log |V| + \alpha(|V|)\right)\right) = \\
        = \left\{ \alpha(|V|) = \O(\log |V|) \right\} = \O(|E| \log |V|)
    \end{multline*}

\end{proof}

\St Результатом функции \texttt{KRUSKAL} является миностов.
\begin{proof}
    
    Утверждается, что после каждой итерации цикла в строчке 6 множество $T$ является подграфом некоторого миностова. 
    Покажем это по индукции.

    База. $T = \varnothing$ -- тривиально.

    Пусть для каждого $0, 1 \dots n - 1 < |V|$ доказали утверждение. 
    Пусть мы рассматриваем на $n$-ом шаге ребро $(u, v)$.
    Если $u$ и $v$ лежат в одной компоненте связности (то есть $\texttt{SAME(}u, v\texttt{)}$ вернул истину), тогда утверждение тривиально, потому что мы ничего не поменяли.
    Пусть они лежали в разных компонентах. Для определённости, $u$ и $v$ лежат в компонентах $A$ и $B$ соответственно.
    Рассмотрим разрез $\mathcal{A} \overset{def}{=} (A, V \setminus A)$. 
    Заметим, что ни одно ребро из $T$ не пересекает $\mathcal{A}$.
    Более того, в силу отсортированности $E$ по весам, ребро $(u, v)$ в частности минимальное ребро, которое пересекает $\mathcal{A}$.
    Также вспомним, что, по предположению индукции, $T$ -- подграф некоторого миностова.
    Значит, по лемме о безопасном ребре, $(u, v)$ -- безопасное ребро для $T \Rightarrow T \cup \{(u, v)\}$ -- подграф некоторого миностова. Доказали требуемое.

    В силу того, что на каждом шаге $T$ -- подграф некоторого миностова, в частности, и после последней итерации $T$ -- подграф некоторого миностова. 
    Но после последней итерации $T$ уже представляет собой дерево, а значит является миностовом.

\end{proof}
